\section{The finite-place interface: completed tests}
\label{sec:interface}

The proposed two-prime interface has now been tested in an explicit
unitary delay model. Its ordinary energy balance and identity limit
hold, but its arithmetic amplitudes fail. A separate argument excludes
a literal finite Euler product with the unchanged archimedean completion.
This section summarizes the derivations in \cite{InterfaceNote}; that
note also contains the detailed path calculations and numerical controls.

\subsection{Primitive-controlled loops and measured amplitudes}

Use the finite Gibbs sector \eqref{eq:finite_hamiltonian} with
$S=\{2,3\}$ and $q_\ell=\ell^{-\beta}$. At each prime attach a
unit-speed transport loop with coordinate $x=\log y$, $1\le y\le\ell$.
Its length and flight time are $\tau_\ell=\log\ell$. This geometry is
an explicit modeling assumption; a thermal energy label alone does not
derive a flight time. In the primitive sector $P_\ell=1$, impose
\begin{equation}
 \binom{y_\ell}{u_\ell}
 =\begin{pmatrix}-r_\ell&t_\ell\\t_\ell&r_\ell\end{pmatrix}
   \binom{f_\ell}{v_\ell},\qquad
 t_\ell=\sqrt{1-r_\ell^2},\quad
 v_\ell(t)=u_\ell(t-\tau_\ell),\quad r_\ell=\ell^{-1}.
 \label{eq:loop_vertex}
\end{equation}
The complementary sector bypasses an empty, disconnected loop. The
full controlled vertex is unitary and preserves the thermal occupations.
Prepare the loops empty, cascade the stages, and measure the mean
outgoing amplitude. The controller is not reset on a repeated traversal.

Writing $z_\ell=e^{-p\tau_\ell}$, elimination of the internal fields gives
\begin{equation}
 F_\ell=\frac{z_\ell-r_\ell}{1-r_\ell z_\ell},\qquad
 R_{\ell,\beta}=q_\ell+(1-q_\ell)F_\ell,\qquad
 R_\beta=R_{2,\beta}R_{3,\beta}.
 \label{eq:thermal_loop_transfer}
\end{equation}
Each $F_\ell$ is a causal contraction, and the product follows from the
independent, unchanged controllers. Set
$d_\ell=q_\ell-r_\ell(1-q_\ell)$ and
$b_\ell=(1-q_\ell)(1-r_\ell^2)$. The coefficients of the paths labeled
$1,2,3,4,6$ are respectively
\begin{equation}
 g_1=d_2d_3,\quad g_2=b_2d_3,\quad g_3=d_2b_3,\quad
 g_4=r_2b_2d_3,\quad g_6=b_2b_3.
\end{equation}
The primitive probability occurs once per distinct prime, but
\begin{equation}
 \frac{g_4}{g_2}=\frac12,\qquad
 \frac{c_4(\beta/2)}{c_2(\beta/2)}=2^{(\beta-1)/2}.
 \label{eq:loop_repetition_failure}
\end{equation}
These ratios differ throughout $0<\beta\le1$. More generally, a fixed
loop reflectivity cannot match the varying ratio on an open temperature
interval. Formally normalizing by $g_1>0$ gives, at $\beta=1$,
\begin{center}
\begin{tabular}{c r r}
\toprule
Label $n$ & $g_n/g_1$ & Required $c_n(1/2)$\\
\midrule
2 & $3/2$ & $1/2$\\
3 & $16/3$ & $2/3$\\
4 & $3/4$ & $1/2$\\
6 & $8$ & $1/3$\\
\bottomrule
\end{tabular}
\end{center}
This normalization compares path weights; it does not define another
contractive readout. Its first-prime shift tangent is $12\log2$,
whereas \eqref{eq:first_prime} requires $(3/2)\log2$. Moreover,
$R_\beta(p)\to d_2d_3>0$ as real $p\to+\infty$, while the completed
target tends to zero like $(2\pi/p)^{\beta/2}$.

\subsection{Energy accounting and a genuine identity limit}

For each controller sector $a$, unitarity of \eqref{eq:loop_vertex}
gives input energy equal to output energy plus loop storage. Let
$y_a$ be the sector output, $\bar y=\mathbb E_\beta y_a=R_\beta f$,
and $E_a(L)$ the total energy remaining in the two loops. Averaging
and separating the output variance yields
\begin{equation}
 \norm{f}_{(0,L)}^2=\norm{\bar y}_{(0,L)}^2
 +\mathbb E_\beta\norm{y_a-\bar y}_{(0,L)}^2
 +\mathbb E_\beta E_a(L).
 \label{eq:loop_energy}
\end{equation}
All terms use ordinary radiation energy. The variance is outgoing
signal fluctuation correlated with the controller, not unexplained
absorption; the controller occupations do not change. This identity
establishes passivity of the specified model, not arithmetic matching.

The model also has the positive endpoint property
\begin{equation}
 \norm{R_\beta-I}\le2\bigl[(1-2^{-\beta})+(1-3^{-\beta})\bigr]
 \longrightarrow0\qquad(\beta\downarrow0).
\end{equation}
Thus a scattering identity limit can exist even when the Gibbs states
have no trace-class limit. This does not repair the amplitude failure.

\subsection{A fixed bounded thermal readout cannot supply centering}

\begin{proposition}[Finite-place readout limitation]
\label{prop:fixed_thermal_readout}
Fix a finite prime set $S$, its native Gibbs state $\rho_\beta$, and
a bounded temperature-independent operator $B$. For any $n>1$
supported on $S$, $\operatorname{Tr}(\rho_\beta B)$ cannot equal
$c_n(\beta/2)$ on a nonempty open interval of positive temperatures.
\end{proposition}
\begin{proof}
In the occupation basis the expectation is
\begin{equation}
 f_B(\beta)=\prod_{\ell\in S}(1-\ell^{-\beta})
 \sum_{\mathbf m\ge0}\langle\mathbf m|B|\mathbf m\rangle
 e^{-\beta\sum_\ell m_\ell\log\ell}.
\end{equation}
Bounded diagonal entries give locally uniform convergence on
$\Ree\beta>0$. Both $f_B$ and $c_n(\beta/2)$ are holomorphic there,
so equality on the interval forces equality throughout this half-plane.
For real $\beta>0$, however, $|f_B(\beta)|\le\norm B$, whereas
$n^{(\beta-1)/2}\prod_{\ell\mid n}(1-\ell^{-\beta})\to\infty$
as $\beta\to\infty$.
\end{proof}
The proposition requires extraction by a fixed bounded observable in
the native Gibbs state. It does not apply automatically to coefficients
obtained by temperature-dependent deconvolution of a completed response,
to a different interacting equilibrium, or to the full $K_{\beta/2}$.

\subsection{Why a finite Euler completion is not passive}

For a finite prime set $S$, put
\begin{equation}
 Z_{S,\beta/2}(p)=\prod_{\ell\in S}
 \frac{1-\ell^{-s_+}}{1-\ell^{-s_-}},\qquad
 d=\frac{1-\beta}{2},\quad s_-=p+d,\quad s_+=p+d+\beta.
\end{equation}
At $p=d$, one has $s_+=1$ and $s_-=1-\beta$. For $0<\beta<1$,
the rational completion in \eqref{eq:archimedean} has a simple pole,
while every finite Euler factor and both gamma factors are finite and
nonzero. Direct evaluation gives
\begin{align}
 \operatorname*{Res}_{p=d}\bigl[A_{\beta/2}Z_{S,\beta/2}\bigr]
 &=-\beta(1-\beta)\pi^{(\beta-1)/2}
 \Gamma\!\left(\frac{1-\beta}{2}\right)
 \prod_{\ell\in S}\frac{1-\ell^{-1}}{1-\ell^{-(1-\beta)}}
 \ne0.
 \label{eq:finite_euler_pole}
\end{align}
A causal time-invariant contraction on ordinary $L^2(0,\infty)$ has
a holomorphic contractive multiplier in $\Ree p>0$. Hence this finite
completion cannot be such a passive response. Additional zero-input
passive ports do not help: their measured scalar block is still
contractive. The pole allows a causal unstable response, so the
exclusion concerns causality together with the required energy bound.

The full quotient cancels the pole through $1/\zeta(s_+)$, using the
simple residue-one pole of zeta at one \cite{DLMF}. In fact
\begin{equation}
 K_{\beta/2}(d)=2\xi(1-\beta).
\end{equation}
This local cancellation is unconditional; it does not establish
full-family passivity. At $\beta=1$, the two-prime completion instead
behaves as $-2/(3\log2\log3\,p^3)$ near $p=0$, whereas
$K_{1/2}(0)=1$. Thus the existing global modular benchmark is unaffected.

Finite \emph{time} windows remain valid comparisons. The far-right
Euler series and the causal archimedean kernel imply equality with the
full formal causal kernel before the first omitted prime delay
\cite[Section~7]{InterfaceNote}. For $S=\{2,3\}$ this means
$L<\log5$. A scalar comparison through $\log6$ must also include
prime 5; a path-resolved two-loop experiment can isolate label 6.
Finite-window agreement alone asserts neither passivity nor a physical
realization.
